\chapter{INTRODUCCIÓN Y PLANTEAMIENTO DEL PROBLEMA}
\label{cap:intro}
\section{Contexto y motivación}
Los cambios de Wikipedia permiten observar cómo una comunidad crea, mantiene y revisa conocimiento. Una edición puede agregar contenido, deshacer una contribución previa o dejar un mensaje en una página de discusión. Al mismo tiempo, los bots realizan tareas automáticas y los registros administrativos describen bloqueos, traslados y agradecimientos. Un conteo total de eventos no explica por sí solo estas diferencias.

El proyecto propone integrar esas señales en una plataforma de consulta. El análisis se limita a \codigo{enwiki} y \codigo{eswiki}, de modo que sea posible comparar dos comunidades y mantener un volumen operativo compatible con los recursos del equipo. El objetivo analítico es describir actividad y relaciones observables; identificar una reversión no equivale a demostrar vandalismo ni a evaluar la calidad de una contribución.

La propuesta responde al Proyecto Grupal II de Big Data: emplear un dataset actualizable de al menos 50\,000 registros, justificar MongoDB, Cassandra y Neo4j, automatizar seis etapas con Airflow, implementar CRUD y presentar indicadores y visualizaciones \cite{enunciado}.

\section{Descripción del problema}
La fuente entrega documentos JSON con estructura variable. Almacenarlos en un único esquema tabular rígido obligaría a tratar numerosos campos opcionales. Consultar repetidamente todos los documentos para dibujar una serie por minuto incrementaría el trabajo de lectura. Por otra parte, preguntas como qué editores comparten páginas o quién revierte a quién se expresan como recorridos entre entidades.

La dificultad comprende también el tiempo. El consumidor recibe eventos durante toda la hora, mientras Airflow consolida ventanas ya cerradas. Una desconexión, un evento reenviado o una corrección manual puede provocar duplicados o resúmenes inconsistentes si ambas capas escriben sin coordinación. El dashboard debe indicar cuándo una cifra es provisional y cuándo procede de una hora consolidada.

La pregunta que guía la implementación es: \textit{¿cómo integrar documentos, series temporales y relaciones para monitorear la actividad colaborativa de Wikipedia con actualización continua y reprocesamiento por ventanas?}

\section{Usuarios y necesidades analíticas}
\begin{table}[H]\centering\small\setstretch{1.05}\renewcommand{\arraystretch}{1.2}
\begin{tabularx}{\textwidth}{P{3.0cm}YY}
\toprule
Usuario & Necesidad & Respuesta del sistema \\
\midrule
Patrulleros y moderadores & Revisar reversiones, cuentas y cambios de contenido. & En vivo, detalle de eventos y red de interacciones. \\
Comunidad editora & Conocer participación, creación y colaboración. & Actividad por editor, páginas nuevas y coedición. \\
Analistas del ecosistema Wikimedia & Comparar actividad humana y automática entre wikis. & Filtro global, automatización y series consolidadas. \\
Operadores del pipeline & Detectar cortes y controlar la calidad de procesamiento. & Estado de las capas, validación, ejecuciones y arquitectura. \\
\bottomrule
\end{tabularx}
\caption{Usuarios previstos y necesidades cubiertas.}
\end{table}

\section{Objetivos}
\subsection{Objetivo general}
Implementar una plataforma multimodelo que capture y procese cambios de Wikipedia, permita consultar su evolución y sus relaciones, y presente indicadores con períodos y fuentes identificables.

\subsection{Objetivos específicos}
\begin{enumerate}
\item Caracterizar una fuente JSON actualizable y documentar su selección, alcance y volumen.
\item Conservar los eventos en MongoDB, materializar consultas temporales en Cassandra y representar interacciones en Neo4j.
\item Automatizar extracción, limpieza, cargas y cálculo de KPIs con un DAG horario de Airflow.
\item Permitir el reprocesamiento de ventanas sin sumar nuevamente sus contribuciones y coordinar la capa rápida con la consolidación.
\item Implementar operaciones CRUD y un dashboard que conecte las tres bases.
\item Definir e interpretar doce indicadores y verificar las reglas de procesamiento y presentación.
\end{enumerate}
